\chapter{Modulation}
\label{ch:modulation}

Any voice can be made a \emph{modulation source}. It then plays like any other voice, from an instrument, through an
envelope and a volume, but it makes no sound: its output steers one parameter of another voice or of the device. A
sample that is a slow sine wave gives vibrato, tremolo or a sweeping filter; a faster or more complex one gives
stranger effects. A modulation source takes a whole voice, so using one leaves five voices for sound.

\section{Modulation Sources}

A voice is a modulation source while its \texttt{MOD\_MODE} bit is 1 (Section~\ref{sec:reg-control}). Its scaled
sample, the same value a sounding voice would send to the mixer (Section~\ref{sec:voice-gain}), is the source value.
The device takes the high eight bits of it, a signed number from $-128$ to $+127$, and calls it $M$. The voice's volume
and envelope therefore set the depth of the modulation, and its sample sets the shape. A voice's \texttt{MOD\_ROUTE}
register says where its value goes; it is read only while \texttt{MOD\_MODE} is 1 (Section~\ref{sec:reg-mod-route}).

\section{Destinations}

\begin{table}[H]
	\centering
	\begin{tabular}{lllp{6.4cm}}
		\toprule
		\textbf{\texttt{DEST}} & \textbf{Destination} & \textbf{Voice field} & \textbf{Parameter steered}                \\
		\midrule
		0                      & Off                  & Ignored              & The voice is a source with no destination \\
		1                      & Pitch                & Used                 & The step of the destination voice         \\
		2                      & Volume               & Used                 & The volume of the destination voice       \\
		3                      & Pan                  & Used                 & The pan of the destination voice          \\
		4                      & Filter cutoff        & Ignored              & \texttt{FILTER\_CUTOFF}                   \\
		5                      & Master volume        & Ignored              & \texttt{MASTER\_VOL}                      \\
		6                      & Filter pan           & Ignored              & The pan byte of \texttt{FILTER\_VOL\_PAN} \\
		7                      & Reserved             & ---                  & Software must not use it                  \\
		\bottomrule
	\end{tabular}
	\caption{Modulation Destinations}
	\label{tab:mod-destinations}
\end{table}

The filter resonance cannot be modulated. A voice may modulate any voice, including itself.

\section{Scaling}
\label{sec:mod-scaling}

The change to the destination is:

\begin{equation*}
	\Delta = M \times 2^{\,\mathrm{SHIFT} + K}
\end{equation*}

\noindent where \texttt{SHIFT} is the signed field of \texttt{MOD\_ROUTE} ($-8$ to $+7$) and $K$ depends on the destination:
$K = 4$ for pitch and for filter cutoff, and $K = 0$ for volume, pan, master volume and filter pan. A negative exponent is an
arithmetic shift right. The \texttt{INV} bit negates $M$ first. The change is added to the value the destination
holds, and the sum is limited to the destination's range: 0 to 255 for volume, pan, master volume, filter cutoff and
filter pan, and 0 to the largest 16.16 value for the pitch step. The register itself is not changed: it keeps the value the
program wrote, and the sum is the value in force.

For volume and pan, a change of the value in force is a change in the same units as the register, so a shift of 0 with
a source of $+127$ moves the volume by 127 steps of the 256. For the cutoff the unit is a step of the cutoff register
(Chapter~\ref{ch:filter}). For pitch the unit is a raw unit of the 16.16 step.

\subsection{The Depth of Pitch Modulation}
\label{sec:mod-pitch-depth}

The pitch change is added to the step in raw units, so its size as a fraction of the pitch depends on the note. A
source of $+127$ with a shift of 0 adds $127 \ll 4 = 2032$ units. Table~\ref{tab:mod-pitch-depth} shows what that is,
in cents, for the 64-byte waveform of Table~\ref{tab:pitch-table}: a slow sine routed to a low note makes a wide swing,
and the same source routed to a high note a barely audible one. A program that wants an equal vibrato across the range
changes the shift as the octave changes.

\begin{table}[H]
	\centering
	\begin{tabular}{rr}
		\toprule
		\textbf{Note} & \textbf{Change for $M = +127$, shift 0 (cents)} \\
		\midrule
		0             & 773                                             \\
		24            & 228                                             \\
		48            & 59.9                                            \\
		72            & 15.2                                            \\
		96            & 3.8                                             \\
		119           & 1.0                                             \\
		\bottomrule
	\end{tabular}
	\caption{The Depth of Pitch Modulation (64-Byte Waveform)}
	\label{tab:mod-pitch-depth}
\end{table}

\section{Rules}

\begin{itemize}
	\item \textbf{One source for each destination.} At most one modulation source affects any destination. If several
	      voices route to the same destination, the one with the lowest voice number wins and the others have no effect. A
	      source with \texttt{DEST} of 0 does not take part.
	\item \textbf{One sample of delay.} A source's value is used by its destination from the next output sample after the
	      source's byte is read. Modulation can therefore be cascaded, and two voices can modulate each other, with no
	      loop that has to settle within one sample: the result is always determined.
	\item \textbf{The envelope comes first.} The envelope of the destination voice is applied before modulation. Modulating a
	      voice's volume never makes sound from a voice whose envelope is inactive, because the gain is the product of the
	      envelope and the volume (Section~\ref{sec:voice-gain}).
	\item \textbf{The programmed value stays.} A destination holds the value the program wrote. Switching the modulation
	      off, by clearing \texttt{MOD\_MODE} on the source or setting \texttt{DEST} to 0, returns the destination to it with
	      no further action.
\end{itemize}

\section{Making a Low-Frequency Oscillator}

A modulation voice is an ordinary voice, so its rate is set by its note and its instrument's pitch table. The lowest
note of a pitch table built for a waveform's real length is 27.5~Hz, which is too fast for vibrato. The program builds
the modulation instrument's pitch table for a cycle that is shorter than the waveform really is, so that its entries
are smaller: a 64-byte sine with a pitch table built as for a 12.8-byte cycle plays at one fifth of the usual rate,
5.5~Hz at note 0. Chapter~\ref{ch:programming} gives an example.
